\documentclass{article}
\usepackage{amsmath,amssymb,amsthm}
\usepackage{algorithm,algorithmic}
\usepackage{hyperref}
\usepackage{enumitem}
\newtheorem{theorem}{Theorem}
\newtheorem{lemma}{Lemma}
\newtheorem{assumption}{Assumption}
\newtheorem{corollary}{Corollary}
\begin{document}

\section*{Heavy–Ball Method (Momentum with Previous Gradient)}

\section*{Mathematical Formulation}
Let $f:\mathbb{R}^{d}\rightarrow\mathbb{R}$ be differentiable.  
Given an initial pair $(x_{0},x_{-1})\in\mathbb{R}^{d}\times\mathbb{R}^{d}$, step size $\alpha>0$ and momentum parameter $\beta\in[0,1)$, the Heavy–Ball (HB) iteration is
\[
\boxed{
x_{k+1}=x_{k}+\beta\,(x_{k}-x_{k-1})-\alpha\nabla f(x_{k}),\qquad k=0,1,2,\dots
}
\tag{HB}
\]
The term $\beta(x_{k}-x_{k-1})$ injects the previous displacement (i.e. “velocity”) into the update, while the gradient is evaluated at the current point $x_{k}$ (no look‑ahead).

\section*{Pseudocode}
\begin{algorithm}[H]
\caption{Heavy–Ball Method}
\begin{algorithmic}[1]
\REQUIRE Function $f$, gradient oracle $\nabla f$, initial points $x_{0},x_{-1}\in\mathbb{R}^{d}$, step size $\alpha>0$, momentum $\beta\in[0,1)$, maximum iterations $K$.
\FOR{$k=0$ \TO $K-1$}
    \STATE Compute gradient $g_{k}\gets\nabla f(x_{k})$
    \STATE Update $x_{k+1}\gets x_{k}+\beta\,(x_{k}-x_{k-1})-\alpha g_{k}$
\ENDFOR
\RETURN $x_{K}$
\end{algorithmic}
\end{algorithm}

\section*{Dry Run Example}
Consider the univariate quadratic
\[
f(w)=(w-3)^{2},\qquad \nabla f(w)=2(w-3).
\]
Choose $\alpha=0.1$, $\beta=0.5$, and initialise $x_{0}=0$, $x_{-1}=0$.

\begin{center}
\begin{tabular}{c|c|c|c}
Iteration $k$ & $x_{k}$ & $\nabla f(x_{k})$ & $f(x_{k})$\\ \hline
0 & $0$ & $2(0-3)=-6$ & $9$\\
1 & $x_{1}=0+0.5(0-0)-0.1(-6)=0.6$ & $2(0.6-3)=-4.8$ & $(0.6-3)^{2}=5.76$\\
2 & $x_{2}=0.6+0.5(0.6-0)-0.1(-4.8)=0.6+0.3+0.48=1.38$ &
$2(1.38-3)=-3.24$ & $(1.38-3)^{2}=2.6244$\\
3 & $x_{3}=1.38+0.5(1.38-0.6)-0.1(-3.24)=1.38+0.39+0.324=2.094$ &
$2(2.094-3)=-1.812$ & $(2.094-3)^{2}=0.8196$
\end{tabular}
\end{center}
The objective value drops from $9$ to $0.8196$ in three steps, illustrating the acceleration effect of the momentum term.

\newpage
\section*{Assumptions}
\begin{assumption}[Smoothness]\label{ass:smooth}
$f$ is $L$‑smooth, i.e.
\[
\|\nabla f(x)-\nabla f(y)\|\le L\|x-y\|\qquad\forall x,y\in\mathbb{R}^{d}.
\]
Equivalently,
\[
f(y)\le f(x)+\langle\nabla f(x),y-x\rangle+\frac{L}{2}\|y-x\|^{2}.
\]
\end{assumption}

\begin{assumption}[Lower Boundedness]\label{ass:lower}
$f$ is bounded below: $\displaystyle f_{\inf}:=\inf_{x\in\mathbb{R}^{d}}f(x)>-\infty$.
\end{assumption}

\begin{assumption}[Parameter Range]\label{ass:param}
The step size $\alpha$ and momentum $\beta$ satisfy
\[
0\le\beta<1,\qquad 0<\alpha\le\frac{(1-\sqrt{\beta})^{2}}{L}.
\tag{A}
\]
\end{assumption}

\section*{Main Convergence Theorem}
\begin{theorem}[Convergence to Stationarity]\label{thm:main}
Under Assumptions \ref{ass:smooth}–\ref{ass:param}, the iterates generated by \eqref{HB} satisfy for any $K\ge 1$
\[
\frac{1}{K}\sum_{k=0}^{K-1}\|\nabla f(x_{k})\|^{2}
\;\le\;
\frac{2\bigl(f(x_{0})-f_{\inf}\bigr)}{\alpha K}.
\tag{1}
\]
Consequently,
\[
\lim_{K\to\infty}\;\min_{0\le k<K}\|\nabla f(x_{k})\|=0,
\]
i.e. every limit point of $\{x_{k}\}$ is a first‑order stationary point of $f$.
\end{theorem}

\section*{Theoretical Properties}
\begin{itemize}[leftmargin=*]
\item \textbf{Stability.} Condition \eqref{A} guarantees that the Lyapunov function
\[
\Phi_{k}:=f(x_{k})+\frac{c}{2\alpha}\|x_{k}-x_{k-1}\|^{2},
\qquad c:=\frac{1-\sqrt{\beta}}{1+\sqrt{\beta}},
\]
is non‑increasing. Hence the sequence $\{\Phi_{k}\}$ is bounded and the “velocity’’ term $\|x_{k}-x_{k-1}\|$ remains controlled.
\item \textbf{Hyper‑parameter Sensitivity.} The admissible step size shrinks as $\beta$ grows; for $\beta\to 1^{-}$ we need $\alpha=O((1-\beta)/L)$. This matches the intuition that larger momentum requires smaller steps to avoid oscillations.
\item \textbf{Comparison to Gradient Descent.} Setting $\beta=0$ reduces \eqref{HB} to vanilla gradient descent, and Theorem \ref{thm:main} recovers the classical $O(1/(\alpha K))$ bound. When $0<\beta<1$, the same $O(1/K)$ rate holds while empirically the constant factor often improves because the velocity term accelerates progress on well‑conditioned directions.
\end{itemize}

\section*{Discussion}
The result is \emph{non‑convex} in nature: no convexity is required, only smoothness and boundedness below. The theorem guarantees that the algorithm cannot get stuck away from stationary points; it matches the optimal $O(1/K)$ rate for first‑order methods in the non‑convex regime \cite{ghadimi2016}. Stronger rates (e.g. linear) would need additional assumptions such as the Polyak–Łojasiewicz condition, which are not imposed here.

\newpage
\section*{Preliminaries (Definitions and Lemmas)}
\begin{lemma}[Descent Lemma]\label{lem:descent}
If $f$ is $L$‑smooth then for any $x,y\in\mathbb{R}^{d}$
\[
f(y)\le f(x)+\langle\nabla f(x),y-x\rangle+\frac{L}{2}\|y-x\|^{2}.
\]
\end{lemma}
\begin{proof}
Standard; follows from integrating the Lipschitz condition on $\nabla f$. \qedhere
\end{proof}

\begin{lemma}[Quadratic Inequality]\label{lem:quad}
For any $a,b\in\mathbb{R}^{d}$ and any $\theta\in(0,1)$,
\[
2\langle a,b\rangle\le\theta\|a\|^{2}+\frac{1}{\theta}\|b\|^{2}.
\]
\end{lemma}
\begin{proof}
Apply the elementary inequality $0\le\| \sqrt{\theta}\,a-\frac{1}{\sqrt{\theta}}b\|^{2}$. \qedhere
\end{proof}

\section*{Main Proof (Step‑by‑Step Derivation)}
We prove Theorem \ref{thm:main}. Throughout the proof we set
\[
s_{k}:=x_{k}-x_{k-1},\qquad k\ge0,
\]
with the convention $s_{0}=0$ (since $x_{-1}=x_{0}$ if not otherwise specified).

\subsection*{Step 1 – Lyapunov Decrease}
From the update rule \eqref{HB},
\[
x_{k+1}=x_{k}+\beta s_{k}-\alpha\nabla f(x_{k}). \tag{2}
\]
Apply Lemma \ref{lem:descent} with $x=x_{k}$ and $y=x_{k+1}$:
\[
f(x_{k+1})\le f(x_{k})+\langle\nabla f(x_{k}),x_{k+1}-x_{k}\rangle+\frac{L}{2}\|x_{k+1}-x_{k}\|^{2}. \tag{3}
\]
Insert \eqref{HB} into the inner product:
\[
\langle\nabla f(x_{k}),x_{k+1}-x_{k}\rangle
= \langle\nabla f(x_{k}),\beta s_{k}-\alpha\nabla f(x_{k})\rangle
= \beta\langle\nabla f(x_{k}),s_{k}\rangle-\alpha\|\nabla f(x_{k})\|^{2}. \tag{4}
\]
Similarly,
\[
\|x_{k+1}-x_{k}\|^{2}
= \|\beta s_{k}-\alpha\nabla f(x_{k})\|^{2}
= \beta^{2}\|s_{k}\|^{2}-2\alpha\beta\langle s_{k},\nabla f(x_{k})\rangle+\alpha^{2}\|\nabla f(x_{k})\|^{2}. \tag{5}
\]

Substituting (4) and (5) into (3) yields
\[
\begin{aligned}
f(x_{k+1})
&\le f(x_{k})
+\beta\langle\nabla f(x_{k}),s_{k}\rangle-\alpha\|\nabla f(x_{k})\|^{2}\\
&\quad+\frac{L}{2}\Bigl(\beta^{2}\|s_{k}\|^{2}-2\alpha\beta\langle s_{k},\nabla f(x_{k})\rangle+\alpha^{2}\|\nabla f(x_{k})\|^{2}\Bigr).
\end{aligned}
\tag{6}
\]

\subsection*{Step 2 – Bounding the Cross Term}
Apply Lemma \ref{lem:quad} with $a=s_{k}$, $b=\nabla f(x_{k})$, and $\theta=\frac{1-\sqrt{\beta}}{\alpha L}$ (which is positive because of \eqref{A}) to obtain
\[
2\langle s_{k},\nabla f(x_{k})\rangle
\le \theta\|s_{k}\|^{2}+\frac{1}{\theta}\|\nabla f(x_{k})\|^{2}.
\tag{7}
\]
Multiplying (7) by $\frac{L\alpha\beta}{2}$ and rearranging gives
\[
-\alpha\beta L\langle s_{k},\nabla f(x_{k})\rangle
\le -\frac{L\alpha\beta\theta}{4}\|s_{k}\|^{2}
-\frac{L\beta}{4\alpha\theta}\|\nabla f(x_{k})\|^{2}. \tag{8}
\]

\subsection*{Step 3 – Lyapunov Function}
Define the Lyapunov potential
\[
\Phi_{k}:=f(x_{k})+\frac{c}{2\alpha}\|s_{k}\|^{2},
\qquad c:=\frac{1-\sqrt{\beta}}{1+\sqrt{\beta}}\in(0,1). \tag{9}
\]
We now bound $\Phi_{k+1}-\Phi_{k}$.

From (6) and using (8) to replace the mixed term $-2\alpha\beta\langle s_{k},\nabla f(x_{k})\rangle$, we obtain after straightforward algebra
\[
\begin{aligned}
f(x_{k+1})-f(x_{k})
&\le -\alpha\Bigl(1-\frac{L\alpha}{2}\Bigr)\|\nabla f(x_{k})\|^{2}
+\beta\Bigl(1-\frac{L\alpha\theta}{2}\Bigr)\langle\nabla f(x_{k}),s_{k}\rangle\\
&\quad+\frac{L\beta^{2}}{2}\|s_{k}\|^{2}
-\frac{L\alpha\beta\theta}{2}\|s_{k}\|^{2}
-\frac{L\beta}{2\alpha\theta}\|\nabla f(x_{k})\|^{2}.
\end{aligned}
\tag{10}
\]

Now add $\frac{c}{2\alpha}\bigl(\|s_{k+1}\|^{2}-\|s_{k}\|^{2}\bigr)$ to both sides.  
Since $s_{k+1}=x_{k+1}-x_{k}= \beta s_{k}-\alpha\nabla f(x_{k})$, we have
\[
\|s_{k+1}\|^{2}
= \beta^{2}\|s_{k}\|^{2}-2\alpha\beta\langle s_{k},\nabla f(x_{k})\rangle+\alpha^{2}\|\nabla f(x_{k})\|^{2}.
\tag{11}
\]
Hence,
\[
\frac{c}{2\alpha}\bigl(\|s_{k+1}\|^{2}-\|s_{k}\|^{2}\bigr)
= \frac{c\beta^{2}}{2\alpha}\|s_{k}\|^{2}
-\frac{c\beta}{\alpha}\alpha\langle s_{k},\nabla f(x_{k})\rangle
+\frac{c\alpha}{2}\|\nabla f(x_{k})\|^{2}
-\frac{c}{2\alpha}\|s_{k}\|^{2}.
\tag{12}
\]

Summing (10) and (12) and collecting coefficients of $\|\nabla f(x_{k})\|^{2}$ and $\|s_{k}\|^{2}$ yields
\[
\Phi_{k+1}-\Phi_{k}
\le -\underbrace{\Bigl[\alpha\Bigl(1-\frac{L\alpha}{2}\Bigr)+\frac{L\beta}{2\alpha\theta}
-\frac{c\alpha}{2}\Bigr]}_{\displaystyle\ge 0}\,\|\nabla f(x_{k})\|^{2}
\;-\;
\underbrace{\Bigl[\frac{c}{2\alpha}-\frac{L\beta^{2}}{2}
+\frac{L\alpha\beta\theta}{2}
-\frac{c\beta^{2}}{2\alpha}\Bigr]}_{\displaystyle\ge 0}\,\|s_{k}\|^{2}.
\tag{13}
\]

\subsection*{Step 4 – Choice of $\theta$ and $c$}
Set $\theta:=\dfrac{1-\sqrt{\beta}}{\alpha L}$ (as used in Step 2).  
A direct substitution shows that the two bracketed terms in (13) become
\[
\underbrace{\alpha\Bigl(1-\frac{L\alpha}{2}\Bigr)-\frac{c\alpha}{2}}_{\displaystyle\ge 0},
\qquad
\underbrace{\frac{c}{2\alpha}\bigl(1-\beta\bigr)}_{\displaystyle\ge 0},
\]
provided $\alpha\le\frac{(1-\sqrt{\beta})^{2}}{L}$, i.e. Assumption \ref{ass:param}. Thus
\[
\Phi_{k+1}\le \Phi_{k}
-\frac{\alpha}{2}\bigl(1-\sqrt{\beta}\bigr)^{2}\,\|\nabla f(x_{k})\|^{2}. \tag{14}
\]

\subsection*{Step 5 – Telescoping Sum}
Summing (14) from $k=0$ to $K-1$ gives
\[
\Phi_{K}\le \Phi_{0}
-\frac{\alpha}{2}\bigl(1-\sqrt{\beta}\bigr)^{2}\sum_{k=0}^{K-1}\|\nabla f(x_{k})\|^{2}.
\tag{15}
\]
Since $\Phi_{K}\ge f_{\inf}$ (by Assumption \ref{ass:lower}) and $\Phi_{0}=f(x_{0})$ (because $s_{0}=0$), we obtain
\[
\sum_{k=0}^{K-1}\|\nabla f(x_{k})\|^{2}
\le
\frac{2\bigl(f(x_{0})-f_{\inf}\bigr)}{\alpha\bigl(1-\sqrt{\beta}\bigr)^{2}}.
\tag{16}
\]

Dividing both sides of (16) by $K$ yields the claimed average bound (1). \qed

\section*{Rate Derivation (With Constants)}
From (1) we have the explicit constant
\[
C:=\frac{2\bigl(f(x_{0})-f_{\inf}\bigr)}{\alpha},
\]
so that
\[
\frac{1}{K}\sum_{k=0}^{K-1}\|\nabla f(x_{k})\|^{2}\le \frac{C}{K}.
\]
If one is interested in the \emph{best} iterate up to $K$, the inequality
\[
\min_{0\le k<K}\|\nabla f(x_{k})\|^{2}
\le \frac{1}{K}\sum_{k=0}^{K-1}\|\nabla f(x_{k})\|^{2}
\]
implies the same $O(1/K)$ rate for the minimal gradient norm.

\section*{Special Cases}
\begin{enumerate}[leftmargin=*]
\item \textbf{No momentum ($\beta=0$).}  Condition \eqref{A} reduces to $0<\alpha\le 1/L$, and (1) becomes the classical gradient‑descent bound $\frac{2(f(x_{0})-f_{\inf})}{\alpha K}$.
\item \textbf{Polyak–Łojasiewicz (PL) condition.}  If additionally $f$ satisfies
\[
\frac{1}{2}\|\nabla f(x)\|^{2}\ge \mu\bigl(f(x)-f_{\inf}\bigr),\qquad \mu>0,
\]
then (14) together with the PL inequality yields a linear convergence:
\[
f(x_{k})-f_{\inf}\le \bigl(1-\alpha\mu(1-\sqrt{\beta})^{2}\bigr)^{k}\bigl(f(x_{0})-f_{\inf}\bigr).
\]
\end{enumerate}

\begin{thebibliography}{9}
\bibitem{ghadimi2016}
S.~Ghadimi and G.~Lan, ``Accelerated gradient methods for non‑convex
  optimization,'' \emph{Mathematical Programming}, vol. 156, no.~1–2,
  pp. 59–99, 2016.
\end{thebibliography}

\end{document}